%! Parametric Functions and Rates of Change — Unit 8, Lesson 8.2 — Homework
\documentclass[10pt]{article}
\usepackage{precalculus-article}
\usepackage{precalculus-boxes}
\newcommand{\TallMath}[1]{$\displaystyle #1\rule[-1.4em]{0pt}{3.2em}$}

\begin{document}

\pageheader{Unit 8, Lesson 8.2}{Homework}
\namedateperiod

\vspace{0.12in}

\begin{notesbox}{Parametric Rates of Change Practice}

\textbf{1.}\quad A particle moves with $h(t) = (t^2 - 4t,\; 2t - t^2 + 1)$
for $0 \le t \le 4$.

\vspace{4pt}
(a)\quad Complete the table.

\vspace{4pt}
\begin{center}
\renewcommand{\arraystretch}{1.6}
\begin{tabular}{|c|c|c|c|}
\hline
$t$ & $x(t) = t^2 - 4t$ & $y(t) = 2t - t^2 + 1$ & $(x,\; y)$ \\
\hline
$0$ & \blank{2cm} & \blank{2cm} & \blank{2.8cm} \\
\hline
$1$ & \blank{2cm} & \blank{2cm} & \blank{2.8cm} \\
\hline
$2$ & \blank{2cm} & \blank{2cm} & \blank{2.8cm} \\
\hline
$3$ & \blank{2cm} & \blank{2cm} & \blank{2.8cm} \\
\hline
$4$ & \blank{2cm} & \blank{2cm} & \blank{2.8cm} \\
\hline
\end{tabular}
\end{center}

\vspace{6pt}
(b)\quad At $t = 0.5$, is the particle moving right or left? Up or down?
Show your reasoning by comparing table values.

\writelines{3}

\vspace{6pt}
(c)\quad At $t = 3$, is the particle moving right or left? Up or down?

\writelines{2}

\end{notesbox}

\vspace{0.08in}

\begin{notesbox}{Problems 2--3: Average Rates of Change and Slope}

\textbf{2.}\quad Use $h(t) = (t^2 - 4t,\; 2t - t^2 + 1)$ from Problem 1.

\vspace{4pt}
(a)\quad Compute the slope of the curve between $t = 0$ and $t = 2$.
Show the computation of $\Delta x$ and $\Delta y$ separately, then
form the ratio.

\writelines{4}

\vspace{6pt}
(b)\quad Compute the slope of the curve between $t = 1$ and $t = 4$.

\writelines{4}

\vspace{0.06in}
\textbf{3.}\quad A different particle moves with
$p(t) = (3t - 6,\; t^2 - 4t + 3)$ for $0 \le t \le 5$.

\vspace{4pt}
(a)\quad Compute the slope of the curve between $t = 1$ and $t = 3$.

\writelines{4}

\vspace{4pt}
(b)\quad At what value of $t$ in $[0, 5]$ is $x(t) = 0$? State the
particle's position at that moment.

\writelines{3}

\end{notesbox}

\vspace{0.08in}

\begin{notesbox}{Problem 4: Two Parametrizations of the Same Line}

\textbf{4.}\quad Consider the straight-line path traced by
$f(t) = (2t,\; 3t - 1)$ for $0 \le t \le 3$.

\vspace{4pt}
(a)\quad State the start point and end point of $f$.

\writeline

\vspace{4pt}
(b)\quad Write a second parametric function $g(t)$ that traces the
\emph{same} straight-line path from the end point back to the start point
as $t$ increases from $0$ to $3$. (Hint: reverse the direction.)

\writelines{3}

\vspace{4pt}
(c)\quad Compute the slope of the curve for $f$ between $t = 0$ and $t = 3$.
Then compute the slope for $g$ between $t = 0$ and $t = 3$ (which are
the reversed endpoints). Are the slopes equal? Explain.

\writelines{4}

\end{notesbox}

\vspace{0.08in}

\begin{notesbox}{Problem 5: AP-Style Multiple Choice}

\textbf{5.}\quad A particle moves in the plane with parametric function
$f(t) = (t^2 - 2t,\; t^2 + 2t)$ for $0 \le t \le 4$.
Which of the following correctly describes the particle's motion at $t = 1$?

\vspace{6pt}
\begin{enumerate}[label=(\Alph*), leftmargin=1.8em, itemsep=4pt]
  \item Moving right and up
  \item Moving left and up
  \item Moving left and down
  \item Moving right and down
\end{enumerate}

\vspace{4pt}
Show your work:

\writelines{3}

\end{notesbox}

\end{document}
